\section{Research Design}

The unit of analysis is a city-platform case rather than a city-year. A case
follows a major LGFV or platform-like local state-owned enterprise through a
documented exit, transformation, restructuring, or termination event and then
locates the relevant function after that event. This unit is narrower than a
city-level exit count, but it makes the measurement problem observable. It can
distinguish a function that ended, remained with the same issuer, moved to
another public entity, or became part of a formal fiscal arrangement.

Sample construction begins with government records, finance-bureau and
state-owned asset supervision documents, exchange disclosures, bond
prospectuses, rating reports, legal opinions, annual reports, business
registration records, and credible financial reporting. A source-only
candidate contains evidence about a platform function but lacks either a formal
event or post-event functional evidence. A label candidate contains both. A
case enters the reference file only after review of the original or
near-original documents under the codebook. This separation prevents source
availability or a single compliance sentence from becoming an outcome by
default.

The current reference file contains ninety-four cases. It was assembled to
develop the classification, expose difficult boundaries, and preserve
documentary variation; it is not a probability sample of Chinese LGFVs.
Guangzhou City Construction Investment and Ningbo Development Investment are
the two substantive exits. Ten cases are functional transfers, and eighty-two
are nominal exits. No liquidation case has yet passed the inclusion rule. The
absence of liquidation and the high nominal-exit share describe the assembled
reference file, not national prevalence.

The explanatory application matches eighty-four reference cases to Ming-Qing
examination-elite density from the CBDB-GADM crosswalk. Elite density is used as
a historical proxy, not as a direct measure of present fiscal administration.
The design reports its association with an indicator for movement away from
nominal exit and examines named cases for contemporary fiscal absorption,
bureaucratic coordination, and project or state-asset governance. These process
codes remain qualitative and cannot establish mediation.

Contemporary controls address the most direct alternative explanations: GDP
per capita, fiscal self-sufficiency, statutory debt pressure, land-finance
dependence, source coverage, platform administrative level, capital or
sub-provincial status, and province. The complete-control subset has
seventy-eight cases and all twelve observed institutional-change outcomes. It
excludes six nominal exits with comparatively low historical capacity. The
initial full-control design also includes two exhaustive platform dummies with
an intercept and is rank deficient. Adjusted results are therefore diagnostic
until the reference category is repaired and rare-outcome sensitivity analyses
are independently reproduced.

The design keeps three inferential objects separate. The ninety-four
human-reviewed cases provide reference labels for measurement development. The
eighty-four historically matched cases provide a nonrepresentative exploratory
association. The LLM screening file expands source discovery but does not
provide observed outcomes. A downstream estimator that combines human and LLM
labels requires a probability validation sample with known inclusion
probabilities and independent human coding. The current post-hoc overlaps do
not meet that requirement.
